\section{``Hebrew Truth'': Philology, Polemic, and Reception}

Jerome's phrase \emph{hebraica veritas} names a methodological and theological claim: for many Old Testament questions, the Hebrew text should govern correction and translation. The phrase marked a consequential reorientation in the Latin West. It did not mean that Jerome possessed an unmediated original, that his every rendering followed Hebrew alone, or that the Septuagint ceased to have ecclesial authority.

\subsection{Fountains, streams, and actual practice}

Jerome often pictures source languages as fountains and translations as streams. The image supports recourse to Hebrew and Greek when Latin witnesses disagree. It is rhetorically powerful but historically incomplete. A translator reads through learned habits, lexica embodied in earlier versions, teachers, parallel passages, and inherited exegesis. Jerome's practice confirms this complexity. Greek Jewish versions could clarify a Hebrew word; the Septuagint could preserve an interpretive tradition; a familiar Latin expression could survive when it conveyed the sense adequately (Kamesar 1993; Rodrigue 2025, especially 309--313). Rodrigue's conclusion is deliberately book-bounded: its direct case studies do not authorize the same mixture or motive in every translated book.

Letter 57's defense of rendering sense for sense, except where even biblical word order may be significant, is similarly contextual. It arose in a dispute about translation and cannot serve as a single algorithm for all his biblical books. Some passages are close and formal; others explain, harmonize, transliterate, or choose among ancient interpretations. Translation technique must be tested book by book.

Jerome's references to Jewish teachers and his use of Jewish Greek translators document real intellectual exchange. His polemical statements about Jews or alleged textual corruption belong to late-antique Christian controversy and cannot be repeated as neutral descriptions of Jewish transmission. The evidence supports contact, borrowing, argument, and asymmetry; it does not disclose every teacher's identity or reconstruct each lesson.

\subsection{Augustine's pastoral and textual objection}

The exchange between Augustine and Jerome gives the dispute a concrete setting. Augustine praised Jerome's ability to revise the Gospels from Greek. He was much more cautious about replacing the Church's Septuagintal Old Testament with a Hebrew-based Latin text. In Letter 71, written in 403, he worried that a novel reading proclaimed in church could disturb congregations and create a gulf between Greek- and Latin-speaking Christians.

The controversy over Jonah's plant made the concern vivid. Augustine's Letter 71 reports that a congregation accustomed to one Latin rendering reacted sharply to Jerome's alternative. Jerome answered in his Letter 112, also transmitted in the Augustinian correspondence as Letter 75, defending his source work and lexical choice. Augustine's Letter 82 then responded to Jerome and restated the pastoral and ecclesial concern while affirming his respect for Jerome's learning. The sequence establishes neither that every congregation reacted in this way nor that one party simply rejected scholarship. It shows two goods in tension: correction from sources and continuity in a worshiping community (Augustine, Letters 71 and 82; Jerome, Letter 112).

Augustine's \emph{City of God} later speaks appreciatively of Jerome as learned in Hebrew, Greek, and Latin while continuing to defend the Septuagint's providential and ecclesial importance. This combination prevents a false binary. One could esteem Jerome's Hebrew translation without making it the sole measure of the Old Testament's Christian reception.

\evidenceframe{Augustine's Letter 71; Jerome's reply, Letter 112 (Augustine, Letter 75); Augustine's further reply, Letter 82; and \emph{City of God} 18.43.}{A prominent contemporary welcomed Jerome's Greek Gospel work but disputed the pastoral and ecclesial consequences of a Hebrew-based Old Testament; the disagreement persisted alongside respect for Jerome's learning.}{The letters must be read in sequence. They report elite correspondence and a narrated congregational incident; they do not quantify Latin Christian opinion or decide the textual question by themselves.}

\subsection{Canon and source are different questions}

Jerome's preference for a Hebrew enumeration of Old Testament books shaped several prefaces, especially the helmeted prologue to Kings. But source-language preference, personal canonical classification, ecclesial reading, and later conciliar reception are distinct historical questions. Augustine's list in \emph{De doctrina christiana} includes books that Jerome treated differently. The later Vulgate transmitted deuterocanonical books, some in Old Latin forms, and Trent received entire books and all their parts as found in the old Latin Vulgate.

It is therefore misleading either to make Jerome the author of the later Catholic canon or to say simply that he ``removed'' its deuterocanonical books. His judgments are primary evidence for one learned late-antique position; the physical and institutional history of the Vulgate records a broader reception.

\subsection{Why the Hebrew project prevailed gradually}

The Hebrew-based books offered learned readers an answer to discrepancies among Latin and Greek copies and increasingly traveled under Jerome's authority. Yet adoption depended on copying networks, schools, monasteries, commentary, and liturgy. The Gallican Psalter's success over the Hebrew Psalter is a particularly strong counterexample to any claim that source priority alone determined reception. A text might be philologically prized, pastorally resisted, liturgically familiar, or materially available in different degrees.

The long-term result was not victory by erasure. It was a Latin Bible in which Hebrew-based translations, Greek-based revisions, and Old Latin inheritances stood side by side. The next chapter makes that composition explicit.
